\newif\ifsolutions
\solutionstrue % Show solutions
%\solutionsfalse % Hide solutions

\documentclass{article}
\usepackage{geometry}
\geometry{margin=1in}
\usepackage{tikz}
\usepackage{amssymb}

% fleqn allows setting indent of display math
\usepackage[fleqn]{amsmath}
\setlength{\mathindent}{0pt} % Set indent
% Disable equation numbering (https://tex.stackexchange.com/a/360378)
\makeatletter
\renewcommand\tagform@[1]{}
\makeatother

% Allow Unicode (some, e.g., © and £ at least)
% https://tex.stackexchange.com/questions/370278/is-there-any-reason-to-use-inputenc
\usepackage[utf8]{inputenc}

% Hyperlinks
\usepackage{hyperref}
\hypersetup{colorlinks=true, urlcolor=blue, linkcolor=blue}

% Prevent indentation of paragraphs
\setlength\parindent{0pt}
\setlength{\parskip}{\baselineskip}

% Spacing above/below equations
% https://tex.stackexchange.com/a/69678
\AtBeginDocument{%
 \abovedisplayskip=-\parskip
 \abovedisplayshortskip=-\parskip
 \belowdisplayskip=0pt
 \belowdisplayshortskip=0pt
}

% Allow 3 additional subsection levels
% https://tex.stackexchange.com/a/60212
\usepackage{titlesec}
\setcounter{secnumdepth}{6}
% H4 in HTML
\titleformat{\paragraph}{\normalfont\normalsize\bfseries}{\theparagraph}{1em}{}
\titlespacing*{\paragraph}{0pt}{3.25ex plus 1ex minus .2ex}{1.5ex plus .2ex}
% H5 in HTML
\titleformat{\subparagraph}{\normalfont\normalsize\bfseries}{\thesubparagraph}{1em}{}
\titlespacing*{\subparagraph}{0pt}{3.25ex plus 1ex minus .2ex}{1.5ex plus .2ex}
% H6 in HTML
\titleformat{\subsubparagraph}{\normalfont\normalsize\bfseries}{\thesubsubparagraph}{1em}{}
\titlespacing*{\subsubparagraph}{0pt}{3.25ex plus 1ex minus .2ex}{1.5ex plus .2ex}

% So enumerate at all levels is numbers
% https://tex.stackexchange.com/questions/78842/nested-enumeration-numbering
\renewcommand{\labelenumii}{\arabic{enumii}.}
\renewcommand{\labelenumiii}{\arabic{enumiii}.}
\renewcommand{\labelenumiv}{\arabic{enumiv}.}

\renewcommand{\mbox}{\text}
\newcommand{\ds}[0]{\displaystyle}
\newcommand{\ihat}[0]{\hat{\boldsymbol{\imath}}}
\newcommand{\jhat}[0]{\hat{\boldsymbol{\jmath}}}
\newcommand{\khat}[0]{\hat{\boldsymbol{k}}}
\newcommand{\xhat}[0]{\hat{\mathbf{x}}}
\newcommand{\yhat}[0]{\hat{\mathbf{y}}}
\newcommand{\zhat}[0]{\hat{\mathbf{z}}}
\newcommand{\rhat}[0]{\hat{\mathbf{r}}}
\newcommand{\bfvec}[1]{\vec{\mathbf{#1}}}
\newcommand{\bfcdot}[0]{\boldsymbol{\cdot}}

\usepackage{fancyhdr}
\pagestyle{fancy}
\lhead{Electric Field}
\rhead{\thepage}
\fancyfoot{}

\begin{document}

% Figures:
% https://www.mathcha.io/editor/M55KMuQLiLmH9Vgp0Ptw6QGB3HnpzJnnuM6VoB7

\section{Overview}

The electric field vector, $\bfvec{E}$, is a quantity assigned to a point in space. Given this quantity, we can compute the force on a charge $Q$ will experience if it is placed at that point using the equation $\bfvec{F}=Q\bfvec{E}$. The direction of $\bfvec{E}$ is also the direction a charge will begin to move if released from rest.

To find $\bfvec{E}$ at any point in space, 
compute the force $\bfvec{F}$ due to all other charges on a hypothetical (or ``test") charge $q$o$ at a point where you want to know $\bfvec{E}$. To find $\bfvec{E}$ at that point, divide $\bfvec{F}$ by $q$o$.

$$\bfvec{E} = \frac{\bfvec{F}}{q$o}$$

%To find $\bfvec{F}$ when a different charge $Q$ is placed where $q$ was, multiply $\bfvec{E}$ by $Q$.

\section{Example I}

Charge $q$1$ is at $(x,y)=(-a,-a)$. Find the electric field at $(x,y)=(a,a)$ in the form $\bfvec{E}=E$x\ihat + E$y\jhat$. Also, find $E$. (Note that $E$ and $|\bfvec{E}|$ are used interchangebly.)

\textbf{Solution}

To find the electric field at a point in space, we put a hypothetical ``test" charge $q$o$ at that point, compute the force on it due to all other charges, and then use

$$\bfvec{E} = \frac{\bfvec{F}}{q$o}$$

The force a positive charge $q$1$ at $(x,y)=(-a,-a)$ exerts on a positive charge $q$2$ at $(x,y)=(a, a)$ was computed in a previous activity.

$$\bfvec{F}${q$1\text{ on } q$2}=k\frac{|q$1q$2|}{8a^2}(\cos 45^\circ \ihat + \sin 45^\circ \jhat)$$

We also found that this equation applies when $q$1$ and $q$2$ are both positive or both are negative. If $q$1$ was positive and $q$2$ was negative, or vice-versa, we found the sign changed:

$$\bfvec{F}${q$1\text{ on } q$2}=-k\frac{|q$1q$2|}{8a^2}(\cos 45^\circ \ihat + \sin 45^\circ \jhat)$$

Based on this, we can write a single equation for all possibilities:

$$\bfvec{F}${q$1\text{ on } q$2}=k\frac{q$1q$2}{8a^2}(\cos 45^\circ \ihat + \sin 45^\circ \jhat)$$

If we replace $q$2$ with $q$o$, this is

$$\bfvec{F}${q$1\text{ on } q$o}=k\frac{q$1q$o}{8a^2}(\cos 45^\circ \ihat + \sin 45^\circ \jhat)$$

The electric field at the location of $q$o$ is then 

$\ds\bfvec{E}${\text{at }(a,a) \text{ due to }q$1} = \frac{\bfvec{F}}{q$o} = \frac{kq$1}{8a^2}(\cos 45^\circ \ihat + \sin 45^\circ \jhat) =\frac{kq$1}{8a^2}\left[\frac{1}{\sqrt{2}}\ihat + \frac{1}{\sqrt{2}}\jhat\right]$,

where the fact that $\sin 45^\circ=\cos 45^\circ=1/\sqrt{2}$ was used. The magnitude is $E=\sqrt{E$x^2+E$y^2}=k|q$1|/8a^2$.

Sign check: When computing electric fields and forces, it is easy to make a sign error. The electric field vector points in the direction a positive charge will move if released there from rest. Suppose $q$1$ is positive. Our equation predicts that a positive charge released from rest at $(a, a)$ will move up and to the right. Suppose $q$1$ is negative. Our equation predicts that a positive charge will move down and to the left. This is consistent with the fact that like charges repel and unlike charges attract.

\section{Problem I}

Charge $q$1$ is at $(x,y)=(-a,a)$. At $(x,y)=(a, 0)$, find $\bfvec{E}$ in the form $\bfvec{E}=E$x\ihat + E$y\jhat$. Check signs of the components of $\bfvec{E}$ using the technique used in Example I. Also, find $E$.



\tikzset{every picture/.style={line width=0.75pt}} %set default line width to 0.75pt        

\begin{tikzpicture}[x=0.75pt,y=0.75pt,yscale=-1,xscale=1]
%uncomment if require: \path (0,208); %set diagram left start at 0, and has height of 208

%Shape: Grid [id:dp33505564538856114] 
\draw  [draw opacity=0] (2,2) -- (552,2) -- (552,202) -- (2,202) -- cycle ; \draw  [color={rgb, 255:red, 220; green, 220; blue, 220 }  ,draw opacity=1 ] (52,2) -- (52,202)(102,2) -- (102,202)(152,2) -- (152,202)(202,2) -- (202,202)(252,2) -- (252,202)(302,2) -- (302,202)(352,2) -- (352,202)(402,2) -- (402,202)(452,2) -- (452,202)(502,2) -- (502,202) ; \draw  [color={rgb, 255:red, 220; green, 220; blue, 220 }  ,draw opacity=1 ] (2,52) -- (552,52)(2,102) -- (552,102)(2,152) -- (552,152) ; \draw  [color={rgb, 255:red, 220; green, 220; blue, 220 }  ,draw opacity=1 ] (2,2) -- (552,2) -- (552,202) -- (2,202) -- cycle ;




\end{tikzpicture}


\ifsolutions
{\bf Answer}:

From Problem II on the Electric Force Activity, if $q$2$ is at $(x,y)=(a, 0)$ (and using the arguments in the previous example to drop the absolute value sign),

$$\bfvec{F}${q$1\text{ on } q$2}=k\frac{q$1q$2}{5a^2}\left(\frac{2}{\sqrt{5}}\ihat -\frac{2}{\sqrt{5}}\jhat\right)$$

Replacing $q$2$ with a test charge $q$o$,

$$\bfvec{F}${q$1\text{ on } q$o}=k\frac{q$1q$o}{5a^2}\left(\frac{2}{\sqrt{5}}\ihat -\frac{2}{\sqrt{5}}\jhat\right)$$

$\ds\bfvec{E}${\text{at }(a,0) \text{ due to }q$1} = \frac{\bfvec{F}}{q$o}=\frac{kq$1}{5a^2}\left(\frac{2}{\sqrt{5}}\ihat -\frac{2}{\sqrt{5}}\jhat\right)$
\fi

\newpage

\section{Problem II}

Charge $q$1$ is at $(x,y)=(-a,a)$. Find the electric field at $(x,y)=(a, 0)$, find $\bfvec{E}$ in the form $\bfvec{E}=E$x\ihat + E$y\jhat$. Check signs of the components of $\bfvec{E}$ using the technique used in Example I. Also, find $E$.



\tikzset{every picture/.style={line width=0.75pt}} %set default line width to 0.75pt        

\begin{tikzpicture}[x=0.75pt,y=0.75pt,yscale=-1,xscale=1]
%uncomment if require: \path (0,208); %set diagram left start at 0, and has height of 208

%Shape: Grid [id:dp33505564538856114] 
\draw  [draw opacity=0] (2,2) -- (552,2) -- (552,202) -- (2,202) -- cycle ; \draw  [color={rgb, 255:red, 220; green, 220; blue, 220 }  ,draw opacity=1 ] (52,2) -- (52,202)(102,2) -- (102,202)(152,2) -- (152,202)(202,2) -- (202,202)(252,2) -- (252,202)(302,2) -- (302,202)(352,2) -- (352,202)(402,2) -- (402,202)(452,2) -- (452,202)(502,2) -- (502,202) ; \draw  [color={rgb, 255:red, 220; green, 220; blue, 220 }  ,draw opacity=1 ] (2,52) -- (552,52)(2,102) -- (552,102)(2,152) -- (552,152) ; \draw  [color={rgb, 255:red, 220; green, 220; blue, 220 }  ,draw opacity=1 ] (2,2) -- (552,2) -- (552,202) -- (2,202) -- cycle ;




\end{tikzpicture}


\ifsolutions
{\bf Answer}:

$\ds\bfvec{E}${\text{at }(a,0) \text{ due to }q$1}=\frac{kq$1}{5a^2}\left(\frac{2}{\sqrt{5}}\ihat -\frac{2}{\sqrt{5}}\jhat\right)$, which matches the solution to Problem I, as expected.
\fi

\end{document}